%!TEX root = ../../thesis.tex

\section{Overview}\label{sec:targets-overview}



% GROUND TRUTH (verified against ~/mugiwara and ~/fuzzyemu/sysroots/resources/cases;
% numbers from figures/generated/tab-campaign-inventory.tex):
%
%   target          | formulation | guest                       | blocks reached | role
%   sequence        | input       | 16-byte passcode case       | ~256 sled      | correctness test
%   cJSON           | corpus      | cJSON_ParseWithLength       | 542--958       | ablation subject
%   libxml2         | corpus      | xmlReadMemory (RECOVER)     | 5085--5419     | transfer
%   picohttpparser  | corpus      | phr_parse_request           | 80--81         | transfer
%
% The three corpus campaigns run four seeds each: cJSON carries 15 arms at 200 learner
% iterations, libxml2 and picohttpparser one arm each at 400
% (\cref{tab:campaign-inventory}). sequence does not appear in that table.

% TODO: Open with the split the chapter is organized around and with the table; the
% per-target sections then only add detail. The table is the most useful single page
% of this chapter --- one row per target, with columns:
%
%   target | formulation | guest / harness | action alphabet | reward | role in evaluation
%
% State the asymmetry immediately rather than letting the reader discover it from the
% section count: one target under the input formulation, three under the corpus
% formulation, and the ablation on one of the latter. That is not an oversight but the
% thesis's scope decision (\cref{sec:methodology-setup}) --- the input formulation is
% carried far enough to establish that the learner works, and the corpus formulation
% is where the fuzzing claims are made. Say it here, once, and do not re-litigate it
% in each section.
%
%   - \cref{sec:targets-input-actions}: sequence is a designed target with a
%     hand-shaped reward. It answers "does our EfficientZero implementation learn a
%     deep, sequentially dependent task at all", and nothing about fuzzing
%     (\cref{sec:target-sequence}). Naming that limit here is the thesis's first
%     defence against the "you designed the reward you then climbed" objection
%     (\cref{sec:discussion-threats}).
%   - \cref{sec:targets-corpus-actions}: three real parsers whose reward is raw
%     coverage and was not designed around the agent. Every fuzzing claim rests on
%     these, and only on these.
%
% Then order the three corpus targets by the thing that actually distinguishes them,
% which is scale rather than designed difficulty: picohttpparser reaches ~81 blocks,
% cJSON ~900, libxml2 ~5300 under the same 16-bucket-per-block instrumentation. That
% two-order-of-magnitude spread is the chapter's real ladder, and it is what makes the
% coverage-granularity limit of \cref{sec:platform-coverage} visible in
% \cref{sec:discussion-threats}.
%
% Say what the three corpus targets SHARE, so it is not repeated three times:
%   - the same corpus MDP (\cref{sec:corpus-actions}) --- place/value phase pair,
%     |A| = 256, 64 mutations = 128 steps per episode, 128 parallel environments,
%   - the same persistent-mode harness shape: boot once, snapshot the guest parked on
%     a fixed-size record read, feed one record per execution, never exit. Written for
%     this thesis, one C file per target under mogi_sysroots.
%   - the same reward (MarginalWithDepth): coverage levels this execution reached that
%     the corpus union had not, plus a weighted bonus per level it out-parses the
%     episode's deepest execution by,
%   - the same coverage instrumentation (\cref{sec:platform-coverage}) and snapshot
%     reset (\cref{sec:platform-snapshots}).
% What differs between them is the guest, the record length and the reward scale ---
% list exactly that, and nothing else.
%
% Close with the role each plays, because it is what makes \cref{chap:results}
% readable: cJSON is the target the configuration is developed on
% (\cref{sec:results-ablations}, 15 arms x 4 seeds), and libxml2 and picohttpparser
% test whether the configuration the ablation selected transfers to a target it was
% never tuned on. State the frozen configuration once here --- constant temperature
% $T = 2$ and root Dirichlet $\xi = 10/|A| = 0.039$ --- and point at
% \cref{tab:campaign-inventory} for the full inventory.
%
% OPEN ITEM: the transferred temperature $T = 2$ is not one of the two constant
% temperatures the cJSON ablation actually ran (1.5 and 2.5), and the two transfer
% campaigns ran 400 iterations against cJSON's 200. Both are honest deviations from
% "the same configuration, unchanged" and must be stated here rather than left for a
% reader to find in \cref{tab:campaign-inventory}.
